% !TEX program = xelatex
% BUSCASAM — Whitepaper
% Compile with: latexmk -xelatex whitepaper.tex   (or: xelatex whitepaper.tex)
% Requires: XeLaTeX, Helvetica Neue installed system-wide (ships with macOS).
\documentclass[12pt,a4paper]{article}

%--------------------------------------------------------------------
% Fonts: Helvetica Neue via fontspec (XeLaTeX) — closest installed match
% to the product's Geist Sans (clean neutral grotesque, like the UI).
%--------------------------------------------------------------------
\usepackage{fontspec}
\setmainfont{Helvetica Neue}
\setsansfont{Helvetica Neue}
\newfontfamily\headingfont{Helvetica Neue}

%--------------------------------------------------------------------
% Language & spacing
%--------------------------------------------------------------------
\usepackage{polyglossia}
\setdefaultlanguage{spanish}
\setotherlanguage{english}

\usepackage{setspace}
\setstretch{1.15}

\usepackage[a4paper,top=2.5cm,bottom=2.5cm,left=2.4cm,right=2.4cm]{geometry}
\usepackage{microtype}

%--------------------------------------------------------------------
% Colour palette
%--------------------------------------------------------------------
% Palette mirrors the product UI tokens (frontend/src/app/globals.css):
% one brand blue + a neutral ramp, no secondary hue.
\usepackage{xcolor}
\definecolor{brand}{HTML}{1D4ED8}      % primary blue — UI --primary
\definecolor{brandlt}{HTML}{1E40AF}    % darker blue — UI --primary-hover
\definecolor{accent}{HTML}{1D4ED8}     % accent folds into the one brand blue
\definecolor{accentlt}{HTML}{DBE6FE}   % pale blue fill — UI --primary-tint-2
\definecolor{tint}{HTML}{EFF4FF}       % faint blue wash — UI --primary-tint
\definecolor{rulec}{HTML}{E5E7EB}      % border — UI --border
\definecolor{rowtint}{HTML}{FAFAFA}    % zebra — UI neutral-50
\definecolor{ink}{HTML}{171717}        % body text — UI --foreground
\definecolor{soft}{HTML}{6B7280}       % muted — UI --muted-foreground

\color{ink}

%--------------------------------------------------------------------
% Tables
%--------------------------------------------------------------------
\usepackage{booktabs}
\usepackage{tabularx}
\usepackage{array}
\usepackage{ragged2e}
\usepackage{colortbl}
\renewcommand{\arraystretch}{1.35}
% small rule separation prevents colortbl row fills from painting over rules
\setlength{\aboverulesep}{2pt}
\setlength{\belowrulesep}{2pt}
\newcolumntype{L}{>{\RaggedRight\arraybackslash}X}
\newcolumntype{C}{>{\Centering\arraybackslash}X}
\newcolumntype{R}{>{\Centering\arraybackslash}X}
\newcolumntype{M}[1]{>{\RaggedRight\arraybackslash}p{#1}}
% Header cell: white bold text on brand fill, with a strut for breathing room.
\newcommand{\thead}[1]{{\color{white}\bfseries\rule[-0.9ex]{0pt}{3.2ex}#1}}
% Strut to give every body row a touch more height.
\newcommand{\bodystrut}{\rule[-0.7ex]{0pt}{2.7ex}}
\arrayrulecolor{rulec}
\setlength{\heavyrulewidth}{1.1pt}
\setlength{\lightrulewidth}{0.6pt}

%--------------------------------------------------------------------
% Section styling
%--------------------------------------------------------------------
\usepackage{titlesec}
\titleformat{\section}
  {\headingfont\Large\bfseries\color{ink}}
  {\color{accent}\thesection\hspace{0.6em}}{0pt}{}
  [\vspace{2pt}{\color{rulec}\titlerule[1pt]}]
\titleformat{\subsection}
  {\headingfont\large\bfseries\color{brandlt}}
  {\thesubsection\hspace{0.5em}}{0pt}{}
\titlespacing*{\section}{0pt}{1.6em}{0.7em}
\titlespacing*{\subsection}{0pt}{1.1em}{0.4em}

%--------------------------------------------------------------------
% Headers / footers
%--------------------------------------------------------------------
\usepackage{fancyhdr}
\pagestyle{fancy}
\fancyhf{}
\renewcommand{\headrulewidth}{0.4pt}
\renewcommand{\headrule}{\hbox to\headwidth{\color{rulec}\leaders\hrule height \headrulewidth\hfill}}
\fancyhead[L]{\footnotesize\headingfont\busmark{0.42pt}\hspace{4pt}\wordmark}
\fancyhead[R]{\footnotesize\color{soft}Whitepaper\,\textbullet\, UNSAM — Grupo 11}
\fancyfoot[C]{\footnotesize\color{soft}\thepage}

%--------------------------------------------------------------------
% Links
%--------------------------------------------------------------------
\usepackage{hyperref}
\hypersetup{
  colorlinks=true,
  linkcolor=brandlt,
  urlcolor=accent,
  citecolor=accent,
  pdftitle={BUSCASAM — Whitepaper},
  pdfauthor={UNSAM — Grupo 11}
}
\usepackage{xurl}
\usepackage{enumitem}
\setlist[itemize]{leftmargin=1.3em,itemsep=2pt,topsep=3pt}
\setlist[enumerate]{leftmargin=1.6em,itemsep=2pt,topsep=3pt}

%--------------------------------------------------------------------
% TikZ (architecture diagram)
%--------------------------------------------------------------------
\usepackage{float}
\usepackage{tikz}
\usetikzlibrary{positioning,arrows.meta,backgrounds,fit,calc,shapes.geometric}

\tikzset{
  box/.style={
    rectangle, rounded corners=3pt, draw=brandlt, line width=0.7pt,
    fill=white, text=ink, align=center, font=\footnotesize,
    inner sep=6pt, minimum height=10mm, text width=26mm},
  store/.style={
    box, fill=accentlt, draw=accent, text width=30mm},
  ext/.style={
    box, draw=soft, dashed, fill=rowtint, text=soft, text width=26mm},
  entry/.style={
    box, fill=brand, text=white, draw=brand, text width=24mm},
  flow/.style={-{Stealth[length=2.2mm]}, line width=0.7pt, draw=brandlt},
  flowa/.style={-{Stealth[length=2.2mm]}, line width=0.7pt, draw=accent},
  biflow/.style={{Stealth[length=2.2mm]}-{Stealth[length=2.2mm]}, line width=0.7pt, draw=brandlt},
  opt/.style={-{Stealth[length=2.2mm]}, line width=0.7pt, draw=soft, dashed},
  lbl/.style={font=\scriptsize\itshape, fill=white, inner sep=1.5pt, text=soft}
}

%--------------------------------------------------------------------
% Brand mark — the magnifying-glass logo from frontend/src/app/icon.svg
% (Wordmark.tsx), redrawn in TikZ. #1 is the per-unit length; the mark
% is 26*#1 tall, matching the 26x26 SVG viewBox (y flipped for TikZ).
%--------------------------------------------------------------------
% baseline: the mark's vertical centre (y=13 units) sits ~0.34em above the
% text baseline, i.e. on the uppercase cap-height midline of the wordmark.
\newcommand{\busmark}[1]{%
  \begin{tikzpicture}[x=#1,y=#1,baseline={13*#1-0.34em}]
    \fill[brand,rounded corners={7.5*#1}] (0,0) rectangle (26,26);
    \draw[white,line width={2.1*#1}] (11,15) circle (4.4);
    \draw[white,line width={2.1*#1},line cap=round] (14.4,11.6) -- (18.5,7.5);
  \end{tikzpicture}%
}
% Two-tone wordmark, matching Wordmark.tsx: BUSCA in ink, SAM in brand blue.
\newcommand{\wordmark}{\textbf{\color{ink}BUSCA\color{brand}SAM}}

%--------------------------------------------------------------------
% Small helpers
%--------------------------------------------------------------------
\newcommand{\lead}[1]{{\large\color{brandlt}#1}}
\usepackage{graphicx}

\begin{document}

%====================================================================
% TITLE PAGE
%====================================================================
\begin{titlepage}
\thispagestyle{empty}
\begin{center}
\vspace*{1.2cm}

{\busmark{2.6pt}}\\[1.1cm]

{\headingfont\fontsize{52}{56}\selectfont\bfseries\color{ink}BUSCA\color{brand}SAM}\\[0.5cm]

{\Large\color{brandlt} Descubrimiento de la producción académica de la UNSAM\\[2pt]
mediante búsqueda híbrida semántica}\\[1.2cm]

{\color{rulec}\rule{0.78\linewidth}{0.7pt}}\\[0.7cm]

{\large\bfseries Trabajo Práctico Final Integrador — Ingeniería de Software}\\[6pt]
{\large Licenciatura en Ciencia de Datos, UNSAM}\\[4pt]
{\large\color{soft} Grupo 11}\\[1.4cm]

\begingroup
\renewcommand{\arraystretch}{1.2}
\setlength{\tabcolsep}{0.45cm}
\begin{tabular}{@{}>{\centering\arraybackslash}p{7.3cm}>{\centering\arraybackslash}p{7.3cm}@{}}
{\bfseries\color{ink}Lucas Achaval}   & {\bfseries\color{ink}Marcos Achaval}\\[-5pt]
{\footnotesize\color{soft}lachavalrodriguez@estudiantes.unsam.edu.ar} &
{\footnotesize\color{soft}machavalrodriguez@unsam-bue.edu.ar}\\[10pt]
{\bfseries\color{ink}Nicolás Cirulli} & {\bfseries\color{ink}Matías Raina}\\[-5pt]
{\footnotesize\color{soft}nlcirulli@unsam-bue.edu.ar} &
{\footnotesize\color{soft}meraina@estudiantes.unsam.edu.ar}\\
\end{tabular}
\endgroup

\vfill
{\color{soft} Junio de 2026}
\end{center}
\end{titlepage}

%====================================================================
% TOC
%====================================================================
\thispagestyle{empty}
{\hypersetup{linkcolor=brand}
\renewcommand{\contentsname}{\headingfont\color{brand}Contenidos}
\tableofcontents}
\clearpage
\setcounter{page}{1}

%====================================================================
% 1. RESUMEN EJECUTIVO
%====================================================================
\section{Resumen ejecutivo}

La producción intelectual de la UNSAM (tesis, papers, trabajos prácticos, apuntes) hoy queda dispersa en repositorios aislados por unidad académica, carpetas compartidas y archivos personales. El conocimiento existe, pero no circula. Los estudiantes no pueden construir sobre lo ya investigado, los docentes no encuentran trabajos previos, y el sector productivo no tiene cómo identificar talento o proyectos con potencial de transferencia.

BUSCASAM es una plataforma web donde la comunidad UNSAM publica y encuentra trabajos académicos. Su diferencial es la experiencia completa, con una publicación sin fricción y una búsqueda que combina significado (semántica) con coincidencia léxica tradicional. En lugar de exigir palabras clave exactas, una consulta en lenguaje natural como \emph{«identificar tipos de flores a partir de sus medidas»} encuentra el trabajo \emph{«Estadística Bayesiana Aplicada — Penguins y Iris»}, que ninguna búsqueda por palabras clave habría conectado con esa consulta.

La primera versión del producto está construida, probada y publicada en \href{https://buscasam.org}{buscasam.org}. Cubre el ciclo completo: autenticación institucional con cuentas Google de la UNSAM, publicación de documentos con extracción automática de texto (incluido OCR para escaneos), sugerencia opcional de resumen y palabras clave con IA, búsqueda híbrida con filtros, coautoría, versionado y moderación.

Los beneficios esperados son una reducción drástica del tiempo de búsqueda de trabajos previos, mayor visibilidad y motivación para los autores, y, en fases posteriores, una vía concreta de vinculación entre la universidad y el sector productivo a través de un \emph{radar de talento} basado en la producción real de la comunidad.

%====================================================================
% 2. CONTEXTO
%====================================================================
\section{Contexto}

\noindent\textbf{Situación actual.}\enspace La UNSAM genera cada año una cantidad significativa de trabajos finales, tesis, papers e informes de cátedra. Esa producción queda confinada a repositorios físicos, sistemas de biblioteca limitados a tesis aprobadas, o directamente en carpetas digitales que no salen de cada materia. No existe un punto único donde un estudiante de una escuela pueda descubrir qué se investigó en otra.

\medskip
\noindent\textbf{Evidencia del problema.}\enspace Tres síntomas concretos:

\begin{enumerate}
  \item \textbf{Invisibilidad:} la mayoría de los trabajos prácticos, apuntes y monografías nunca llega a ningún sistema; el repositorio institucional concentra producción formal (tesis y publicaciones) y deja fuera el grueso de la producción de grado.
  \item \textbf{Fricción de acceso:} los buscadores por palabras clave exigen saber de antemano las palabras exactas. El repositorio institucional de la UNSAM (RI-UNSAM, \href{https://ri.unsam.edu.ar}{\nolinkurl{ri.unsam.edu.ar}}) lo ilustra: una consulta en lenguaje natural como \emph{«identificar tipos de flores a partir de sus medidas»} devuelve resultados sin relación con lo buscado.
  \item \textbf{Desconexión con el entorno:} sin una vitrina unificada, una empresa que busca perfiles o proyectos en un área específica no tiene forma de encontrarlos.
\end{enumerate}

\medskip
\noindent\textbf{Impacto de no resolverlo.}\enspace Se repiten temas de investigación por desconocimiento de trabajos previos, el conocimiento no se acumula de una generación de estudiantes a la siguiente, y la universidad desaprovecha oportunidades de convenios, pasantías y financiamiento que dependen de que su producción sea visible.

\medskip
\noindent\textbf{Cómo el proyecto mejora la situación.}\enspace BUSCASAM centraliza la publicación con una fricción mínima (el autor sube un archivo y el sistema hace el resto, extrayendo el texto, derivando el año y sugiriendo resumen y palabras clave) y resuelve el problema de descubrimiento con búsqueda por significado. Al bajar el costo de publicar y de encontrar, convierte producción dormida en una red de conocimiento viva.

%====================================================================
% 3. PROPUESTA DE SOLUCIÓN
%====================================================================
\section{Propuesta de solución}

\noindent\textbf{Propósito.}\enspace Que cualquier integrante de la comunidad UNSAM pueda publicar su trabajo en minutos y que cualquier persona pueda encontrarlo preguntando en lenguaje natural, con garantías de control de acceso y con moderación del catálogo.

\medskip
\noindent\textbf{Funcionalidades principales} (todas operativas en producción):

\begin{itemize}
  \item \textbf{Búsqueda híbrida:} una única caja de texto libre, con filtros opcionales por área de estudio, tipo de documento y fecha. El ranking fusiona similitud semántica y coincidencia léxica. Si el componente semántico no responde, la búsqueda léxica sigue funcionando sola y el buscador permanece disponible.
  \item \textbf{Publicación asistida:} el autor crea un borrador, sube un archivo, y el sistema procesa el documento en segundo plano, con extracción de texto, OCR si es un escaneo, derivación del año y, solo si el autor lo activa, sugerencia de resumen y palabras clave con IA. El autor siempre revisa y edita antes de publicar; nada aparece en la búsqueda sin su confirmación.
  \item \textbf{Control de acceso por niveles:} cada documento es público (cualquier visitante), interno (comunidad UNSAM autenticada) o privado (solo autores). Ninguna respuesta del sistema (búsqueda, conteos, relacionados, sitemap) revela documentos que el visitante no está autorizado a ver.
  \item \textbf{Coautoría y versiones:} invitación de coautores con aceptación explícita, y reemplazo versionado de archivos donde la versión publicada anterior sigue disponible hasta que la nueva esté procesada y confirmada.
  \item \textbf{Moderación:} cualquier usuario autenticado puede reportar un documento; los docentes pueden ocultarlo o restituirlo, con registro de quién, por qué y cuándo.
  \item \textbf{Descubrimiento:} portada con los trabajos más leídos de la semana y “trabajos relacionados” en cada detalle.
\end{itemize}

\medskip
\noindent\textbf{Público objetivo.}

\begin{itemize}
  \item \textbf{Usuarios finales:} estudiantes (publican y buscan), docentes e investigadores (publican, buscan, moderan) y visitantes externos (buscan y descargan lo público).
  \item \textbf{Beneficiarios institucionales:} bibliotecas y secretarías académicas (menos tiempo en gestión de antecedentes), y el rectorado (visibilidad de la producción de la universidad).
  \item \textbf{Socios clave:} las unidades académicas, que aportan el contenido inicial y la difusión; y, en fases futuras, empresas y organismos interesados en identificar talento y proyectos transferibles.
  \item \textbf{Canales:} aplicación web pública (buscasam.org), acceso institucional con las cuentas Google que la comunidad ya usa a diario, y difusión a través de cátedras y bibliotecas.
\end{itemize}

%====================================================================
% 4. IMPACTO ESPERADO Y BENEFICIOS
%====================================================================
\section{Impacto esperado y beneficios}

\noindent\textbf{Comparación con soluciones existentes.}\enspace Los repositorios institucionales cubren solo tesis aprobadas, con carga centralizada y búsqueda por campos exactos. Google Scholar indexa publicaciones formales, no la producción de grado. Las carpetas compartidas por materia no son buscables ni públicas. BUSCASAM se distingue en tres ejes. Cubre toda la producción (del apunte a la tesis), la publica el propio autor en minutos, y busca por significado y no solo por coincidencia textual, como ilustra RI-UNSAM, donde esa misma consulta en lenguaje natural no devuelve el trabajo relevante y BUSCASAM sí.

\medskip
\noindent\textbf{Beneficios esperados:}
\begin{itemize}
  \item \textbf{Reducción sustancial del tiempo de búsqueda de trabajos previos} (a medir en el piloto), al reemplazar la búsqueda fragmentada por palabras clave en catálogos y carpetas dispersas por una única consulta en lenguaje natural con filtros.
  \item \textbf{Aumento de convenios de pasantías y vinculación}, una vez habilitada la fase de vinculación con el sector productivo.
  \item \textbf{Eficiencia administrativa:} menos horas de personal bibliotecario y docente dedicadas a localizar y verificar antecedentes.
\end{itemize}

\medskip
\noindent\textbf{Beneficios cualitativos:} reputación institucional (una universidad que hace circular y accesible su propia producción), sentido de pertenencia y motivación de los estudiantes al ver su trabajo publicado y leído, y mejora de la calidad académica al reducir la redundancia de temas.

\medskip
\noindent\textbf{Indicadores de éxito.}\enspace El Anexo~D define cinco indicadores con su forma de cálculo: trabajos publicados por unidad académica (crecimiento sostenido, mensual), tasa de éxito de búsqueda o proporción de búsquedas con clic en un resultado ($>$60\%, mensual), latencia de búsqueda en el percentil 95 ($<$2 segundos, continua), lecturas semanales de documentos (crecimiento sostenido) y satisfacción de usuarios (encuesta semestral, $\geq$70\% positiva).

%====================================================================
% 5. ARQUITECTURA Y DISEÑO
%====================================================================
\section{Arquitectura y diseño}

\noindent\textbf{Diseño de alto nivel.}\enspace El sistema tiene cuatro piezas (diagrama en el Anexo~A):

\begin{enumerate}
  \item \textbf{Aplicación web} (Next.js/React): búsqueda, detalle, publicación, panel del autor y moderación.
  \item \textbf{API} (FastAPI/Python): toda la lógica de negocio, autenticación y control de acceso.
  \item \textbf{PostgreSQL} como única base de datos, que cumple a la vez cuatro roles, el de almacén de datos transaccionales, el de índice de búsqueda full-text en español, el de índice vectorial de embeddings (extensión pgvector) y el de cola de trabajos en segundo plano (extracción de texto, OCR, indexado).
  \item \textbf{Servicio de embeddings} (modelo multilingüe servido localmente en CPU) que convierte textos en vectores para la búsqueda semántica.
\end{enumerate}

Todo corre en una única máquina virtual en Google Cloud, detrás de un balanceador con HTTPS administrado. La infraestructura completa está definida como código y el despliegue es automático. Cada versión etiquetada se valida (tests, auditoría de dependencias, migraciones), se construye y se publica en producción sin intervención manual.

\medskip
\noindent\textbf{Decisiones de diseño y su justificación.}\enspace El criterio rector fue \emph{máxima confiabilidad con mínima superficie operativa}, porque el equipo es de cuatro personas a tiempo parcial y el presupuesto tiende a cero.

\begin{itemize}
  \item \textbf{PostgreSQL para todo (datos, búsqueda, vectores, colas)} en lugar de sumar un motor de búsqueda, una base vectorial y un broker dedicados. Cada sistema extra es algo más que instalar, monitorear, respaldar y que puede fallar. A la escala actual y proyectada, PostgreSQL resuelve los cuatro problemas con margen, y consolidarlo vuelve el respaldo diario y la restauración una sola operación verificada.
  \item \textbf{Búsqueda híbrida con anclaje léxico.} La búsqueda puramente semántica le asigna similitud a casi cualquier cosa, y una consulta fuera de tema devuelve resultados que “suenan parecido” pero no sirven. Por eso un resultado solo aparece si hay coincidencia léxica, o si supera un piso semántico calibrado \emph{y además} comparte señal textual con la consulta. Preferimos no mostrar nada, y decirlo con claridad, antes que resultados irrelevantes; la confianza en el buscador es el activo central del producto.
  \item \textbf{Procesamiento en segundo plano.} La extracción de texto y el OCR pueden tardar minutos en un PDF escaneado. La API responde de inmediato y procesa después, con el estado visible en el tablero del autor, para que nadie quede mirando una pantalla colgada.
  \item \textbf{IA opcional y por documento.} La sugerencia de resumen y palabras clave es una decisión del autor en cada borrador. Si no la activa, el texto de su documento no sale del sistema. Hay trabajos con datos sensibles, y la confianza de los autores vale más que la comodidad por defecto.
  \item \textbf{Publicar primero, moderar después.} Exigir aprobación previa habría frenado la adopción inicial. El riesgo de contenido inadecuado se mitiga con reporte de usuarios, ocultamiento inmediato por docentes y registro auditable de cada acción; como recaudo, la salida a producción se condicionó a que moderación y borrado estuvieran operativos.
  \item \textbf{Una sola VM en lugar de servicios distribuidos.} A este volumen, distribuir es pagar complejidad sin recibir nada a cambio. La arquitectura interna ya está separada por módulos (búsqueda, publicación, acceso, moderación), así que escalar a más máquinas, si el crecimiento lo exige, es un cambio de infraestructura y no una reescritura.
\end{itemize}

\medskip
\noindent\textbf{Atributos de calidad, en orden de prioridad:}
\begin{enumerate}
  \item \textbf{Seguridad y control de acceso.} Es el atributo no negociable, porque un solo documento privado expuesto destruye la confianza de todos los autores. Cada vía de acceso (búsqueda, detalle, relacionados, descarga, conteos, sitemap) aplica la misma política de visibilidad, y los criterios de aceptación se escribieron en esos términos.
  \item \textbf{Idoneidad funcional de la búsqueda.} El producto es el buscador; si los resultados decepcionan, no hay segunda visita. La función central, devolver lo relevante para cada consulta, se valida con juegos de consultas de prueba con resultados esperados: la relevancia es medible y reproducible, no una impresión.
  \item \textbf{Confiabilidad y tolerancia a fallos.} La búsqueda degrada de forma controlada. Si el componente semántico no responde, la búsqueda léxica sigue operativa. Además, publicar o reemplazar nunca expone una versión sin indexar y una versión candidata fallida no oculta la publicada, de modo que el catálogo permanece consistente ante fallas parciales.
  \item \textbf{Rendimiento.} Respuesta de búsqueda bajo 2 segundos en el percentil 95, a la escala actual y proyectada del catálogo. Un buscador lento no se usa, mientras que una funcionalidad faltante se agrega después.
  \item \textbf{Usabilidad.} Una sola caja de búsqueda, mensajes claros cuando no hay resultados y publicación guiada paso a paso. Las anteriores son condiciones de existencia del producto; esta es condición de crecimiento.
  \item \textbf{Mantenibilidad.} Módulos con límites claros, decisiones de arquitectura registradas por escrito y una suite de pruebas automatizadas (unitarias, de integración y de punta a punta en navegador) que corre en cada cambio.
\end{enumerate}

\medskip
\noindent\textbf{Estrategia de implementación.}\enspace Producto mínimo en producción primero, visión completa por fases. La fase 1 (publicación + búsqueda + moderación, hoy operativa en buscasam.org) valida la hipótesis central (que la búsqueda semántica sobre producción real cambia el descubrimiento) con usuarios reales y costo mínimo. Las fases siguientes (recomendaciones personalizadas, grafos de conocimiento, radar de talento para el sector productivo) se construyen sobre esa base solo si las métricas de adopción de la fase 1 lo justifican. Invertir en la visión completa antes de validar la base habría sido apostar todo el esfuerzo a supuestos sin probar.

%====================================================================
% 6. PLAN DE DESARROLLO
%====================================================================
\section{Plan de desarrollo y cronograma}

\noindent\textbf{Organización y método de trabajo.}\enspace Equipo de cuatro integrantes a tiempo parcial cubriendo seis roles (líder/arquitectura, backend, frontend, IA/datos, infraestructura, calidad). El proyecto abarcó unas $\sim$9 semanas de la materia, pero gracias a la asistencia intensiva de IA agéntica la construcción en código se concentró en unas dos semanas, comprimiendo una carga estimada de $\sim$67 días-persona. Se trabajó en ciclos con entregas verticales, donde cada incremento cruza interfaz, API y datos y termina integrado, probado y desplegable, con revisión de código y pruebas automáticas antes de integrarse.

\medskip
\noindent\textbf{Fases, hitos y cronograma.}\enspace El trabajo se organizó en nueve fases (F0 a F8) sobre el calendario de la materia, desde la definición de alcance y arquitectura hasta la salida a producción, pasando por acceso, búsqueda, publicación, detalle, colaboración, ciclo de vida y descubrimiento, cada una con un hito verificable (detalle en el Anexo~E). La ruta crítica fue acceso $\rightarrow$ búsqueda $\rightarrow$ publicación $\rightarrow$ moderación $\rightarrow$ producción, ya que la búsqueda con visibilidad exige el sistema de roles y la salida a producción se condicionó a que moderación y borrado estuvieran operativos por requisitos de privacidad.

\medskip
\noindent\textbf{Costos.}\enspace El desarrollo se absorbió como trabajo académico del equipo; el único desembolso directo fue la suscripción a la herramienta de IA agéntica. El costo recurrente de operación, estimado sobre la infraestructura desplegada, es:

\medskip
\noindent\begin{tabularx}{\linewidth}{L >{\Centering\arraybackslash}p{5cm}}
\rowcolor{brand}
\thead{Concepto} & \thead{Estimación mensual (USD)} \\
Máquina virtual (2 vCPU, 8 GB) & $\sim$50 \\
\rowcolor{rowtint}
Disco persistente (50 GB) & $\sim$5 \\
Balanceador HTTPS & $\sim$20 \\
\rowcolor{rowtint}
Red, registro de imágenes, secrets, backups & $\sim$8 \\
IA para resúmenes (opcional, por documento) & $<$\,5 \\
\rowcolor{rowtint}
Dominio (prorrateado) & $\sim$1 \\
\midrule
\rowcolor{accentlt}
{\bfseries\bodystrut Total} & \textbf{$\sim$90} \\
\bottomrule[\heavyrulewidth]
\end{tabularx}

\medskip
El servicio de embeddings corre en la propia VM con un modelo abierto, sin costo por consulta, para que el costo operativo no escale con el uso de la búsqueda. El único componente de IA con costo por uso es la sugerencia de resumen y palabras clave, hoy provista por Gemini (Vertex AI), más barata al volumen actual que mantener un modelo propio siempre encendido. El conector para servir ese mismo modelo localmente con pesos abiertos ya está desarrollado, así que trasladar esa carga a infraestructura propia, si el volumen creciera, sería un cambio de configuración, sin código nuevo.

\medskip
\noindent\textbf{Recursos necesarios.}\enspace \emph{Humanos:} el equipo de cuatro para la evolución del producto, más $\sim$8 horas mensuales de un responsable técnico para respaldos, actualizaciones de seguridad y soporte básico; al ser de código abierto, podrían sumarse contribuciones de la comunidad, no garantizadas. \emph{Tecnológicos:} la infraestructura cloud descripta y herramientas de desarrollo asistido por IA.

\medskip
\noindent\textbf{Riesgos principales y respuesta} (registro completo en el Anexo~C). El riesgo dominante es un catálogo inicial pobre que derive en búsquedas decepcionantes y baja adopción; se mitiga con contenido semilla antes de difundir, campaña con cátedras en el lanzamiento y “relacionados” cuando hay pocos resultados. La relevancia del ranking se sostiene con juegos de consultas de prueba, calibración del piso semántico y la búsqueda léxica siempre disponible como respaldo. La inexperiencia en búsqueda semántica se atacó con una prueba de concepto temprana de embeddings e indexado y margen en las estimaciones. Los PDFs escaneados o malformados se cubren con OCR y procesamiento reintentable, con revisión del autor antes de publicar. Por último, el contenido inadecuado por la publicación inmediata se contiene con reporte de usuarios, ocultamiento docente y moderación operativa como condición de lanzamiento.

%====================================================================
% 7. CONCLUSIONES
%====================================================================
\section{Conclusiones y recomendaciones}

La UNSAM produce conocimiento que hoy no circula, y el costo de esa invisibilidad se paga en investigación redundante, oportunidades de vinculación perdidas y trabajo estudiantil que muere en una carpeta. BUSCASAM ataca el problema en su raíz (el costo de publicar y el costo de encontrar) y lo hace con una primera versión funcional, probada y operativa en \href{https://buscasam.org}{buscasam.org}, construida por un equipo de cuatro personas y con un costo de operación del orden de USD 90 mensuales.

\medskip
\noindent\textbf{Próximos pasos sugeridos:}
\begin{enumerate}
  \item \textbf{Piloto formal en una unidad académica:} carga semilla de trabajos existentes y difusión con cátedras, midiendo adopción y tasa de éxito de búsqueda durante un cuatrimestre.
  \item \textbf{Decisión de fase 2 basada en métricas:} si la adopción acompaña, incorporar recomendaciones personalizadas, historial de búsqueda y favoritos (hoy fuera de alcance por diseño).
  \item \textbf{Fase de vinculación:} habilitar la vitrina hacia el sector productivo (perfiles de autor, radar de talento, contacto mediado), que es donde se materializan los beneficios económicos para la institución.
  \item \textbf{Institucionalización:} acordar con bibliotecas y secretarías la integración del depósito de tesis al flujo de BUSCASAM, y formalizar el responsable técnico de operación.
\end{enumerate}

%====================================================================
% 8. ANEXOS
%====================================================================
\clearpage
\section{Anexos}

\subsection*{Anexo A — Arquitectura del sistema}
\addcontentsline{toc}{subsection}{Anexo A — Arquitectura del sistema}

\begin{figure}[H]
\centering
\resizebox{\ifdim\width>\linewidth \linewidth\else \width\fi}{!}{%
\begin{tikzpicture}[node distance=8mm and 10mm]

  % --- Edge / entry (outside the VM) ---
  \node[entry] (user) {Usuario};
  \node[box, below=7mm of user] (lb) {Balanceador GCP};

  % --- Inside the VM ---
  \node[box, below=11mm of lb] (nginx) {nginx};
  \node[box, below left=9mm and 3mm of nginx] (fe) {Web\\Next.js};
  \node[box, below right=9mm and 3mm of nginx] (api) {API\\FastAPI};
  \node[store, below=9mm of fe] (blob) {Almacenamiento};
  \node[store, below=9mm of api] (pg) {PostgreSQL};
  \coordinate (mid) at ($(blob)!0.5!(pg)$);
  \node[box, below=10mm of mid] (workers) {Workers de procesamiento};
  \node[store, below=8mm of workers] (tei) {Servicio de embeddings};

  % --- External services ---
  \node[ext, left=12mm of fe] (ga) {Google OAuth};
  \node[ext] (vrt) at (ga |- tei) {Vertex AI};

  % --- VM container box (behind the inner nodes) ---
  \begin{scope}[on background layer]
    \node[draw=brand, line width=1pt, rounded corners=5pt, dashed,
          fill=brand!4, inner xsep=6mm, inner ysep=5mm,
          fit=(nginx)(fe)(api)(blob)(pg)(workers)(tei)] (vm) {};
    \node[font=\footnotesize\bfseries\itshape, text=brand, anchor=north west]
          at ([xshift=3mm,yshift=-2pt]vm.north west) {Máquina virtual GCP};
  \end{scope}

  % --- Edges ---
  \draw[flow] (user) -- node[lbl]{HTTPS} (lb);
  \draw[flow] (lb) -- (nginx);
  \draw[flow] (nginx.south) -- (fe.north);
  \draw[flow] (nginx.south) -- (api.north);
  \draw[flow] (fe) -- node[lbl]{SSR} (api);

  \draw[flow] (api) -- (pg);
  \draw[flow] (api) -- (blob);
  \draw[flow] (workers) -- (blob);
  \draw[flow] (workers) -- (pg);
  \draw[flowa] (workers) -- (tei);
  \draw[flowa] (api.east) to[bend left=58] (tei.east);

  \draw[opt] (workers.west) to[bend right=12] (vrt.east);
  \draw[biflow] (ga.south) to[bend right=17] (api.west);

  % Phantom point mirroring the left-hand external boxes (Google OAuth / Vertex AI)
  % across the central spine, so the figure's bounding box is symmetric about the
  % nginx/workers/embeddings column and that column centers on the page.
  \node[inner sep=0pt, minimum size=0pt] at ($(ga.west)!2!(ga.west -| nginx)$) {};

\end{tikzpicture}}
\end{figure}

\noindent\textbf{Despliegue:} infraestructura definida como código (Terraform), imágenes construidas y publicadas por CI en cada versión etiquetada, despliegue automático con migraciones y verificación de salud. Respaldo diario con puntos de recuperación apareados (base de datos + archivos) y restauración ensayada antes del lanzamiento.

\subsection*{Anexo B — Plan de pruebas (resumen)}
\addcontentsline{toc}{subsection}{Anexo B — Plan de pruebas}

\begin{itemize}
  \item \textbf{Niveles:} pruebas unitarias (lógica de ranking, validaciones, etapas de indexado), de integración (endpoints contra base de datos real, trabajos en segundo plano, migraciones en ambos sentidos) y de punta a punta en navegador (flujos completos de búsqueda, publicación, coautoría y moderación, incluyendo móvil).
  \item \textbf{No funcionales:} carga concurrente sobre la búsqueda, control de visibilidad por rol en todas las superficies, recuperación ante falla de una etapa de procesamiento, compatibilidad de formatos de archivo.
  \item \textbf{Puertas de calidad:} ningún cambio se integra sin revisión de pares y suite completa en verde; ninguna versión sale a producción sin el conjunto completo de pruebas, auditoría de dependencias y validación de migraciones.
  \item \textbf{Criterios de aceptación clave:} un visitante nunca observa documentos que no le corresponden por ninguna vía; una versión candidata fallida nunca oculta la versión publicada; sin servicio semántico, la búsqueda léxica sigue disponible.
\end{itemize}

\subsection*{Anexo C — Registro de riesgos}
\addcontentsline{toc}{subsection}{Anexo C — Registro de riesgos}

\medskip
\noindent\begingroup
\setlength{\tabcolsep}{6pt}
\renewcommand{\arraystretch}{1.5}
% Whole words instead of mid-word hyphenation in the narrow columns.
\hyphenpenalty=10000\exhyphenpenalty=10000
\begin{tabularx}{\linewidth}{>{\centering\arraybackslash\bfseries}p{0.6cm} >{\RaggedRight\arraybackslash}p{1.85cm} L >{\Centering\arraybackslash}p{1.2cm} >{\Centering\arraybackslash}p{1.75cm} >{\RaggedRight\arraybackslash}p{2.6cm}}
\rowcolor{brand}
\thead{\#} & \thead{Tipo} & \thead{Riesgo} & \thead{Prob.} & \thead{Impacto} & \thead{Responsable} \\
R1 & Negocio & \bodystrut Contenido inicial insuficiente $\rightarrow$ resultados pobres $\rightarrow$ baja adopción & Alta & Alto & Producto \\
\rowcolor{rowtint}
R2 & Técnico & \bodystrut Ranking híbrido percibido como poco relevante & Media & Alto & Búsqueda/IA \\
R3 & Proyecto & \bodystrut Inexperiencia en embeddings e indexado $\rightarrow$ atrasos & Alta & Medio & Líder técnico \\
\rowcolor{rowtint}
R4 & Técnico & \bodystrut PDFs escaneados/malformados $\rightarrow$ extracción incompleta & Media & Medio & Ingesta \\
R5 & Negocio & \bodystrut Publicación inmediata $\rightarrow$ contenido inadecuado visible antes de moderar & Media & Medio & Producto \\
\rowcolor{rowtint}
R6 & Técnico & \bodystrut Datos de uso para recomendaciones $\rightarrow$ riesgos de privacidad & Baja & Alto & Líder técnico \\
R7 & Técnico & \bodystrut Degradación de tiempos de búsqueda con catálogo grande & Baja & Medio & Búsqueda/IA \\
\bottomrule[\heavyrulewidth]
\end{tabularx}\endgroup

\medskip
Los planes de mitigación de R1–R5 están resumidos en la sección 6. R6 se difirió porque la fase 1 no almacena historial de búsqueda personal. R7 se monitorea con el indicador de latencia p95.

\subsection*{Anexo D — Indicadores y forma de cálculo}
\addcontentsline{toc}{subsection}{Anexo D — Indicadores y forma de cálculo}

\medskip
\noindent\begin{tabularx}{\linewidth}{>{\bfseries\RaggedRight\arraybackslash}p{4.5cm} L}
\rowcolor{brand}
\thead{Indicador} & \thead{Cálculo} \\
Trabajos publicados por unidad & \bodystrut Conteo de documentos publicados por mes, segmentado por escuela/área \\
\rowcolor{rowtint}
Tasa de éxito de búsqueda & \bodystrut Búsquedas con al menos un clic en resultado $\div$ búsquedas totales (analítica de producto, agregada y anónima) \\
Latencia p95 de búsqueda & \bodystrut Percentil 95 del tiempo de respuesta del endpoint de búsqueda, medido en producción \\
\rowcolor{rowtint}
Lecturas semanales & \bodystrut Aperturas de detalle deduplicadas (a lo sumo una por lector, por documento, por día), ventana móvil de 7 días \\
Satisfacción & \bodystrut Encuesta semestral a estudiantes y docentes (escala 1–5; meta: $\geq$70\% en 4–5) \\
\bottomrule[\heavyrulewidth]
\end{tabularx}

\subsection*{Anexo E — Fases, hitos y cronograma}
\addcontentsline{toc}{subsection}{Anexo E — Fases, hitos y cronograma}

\medskip
\noindent\begin{tabularx}{\linewidth}{>{\bfseries\RaggedRight\arraybackslash}p{2.6cm} L >{\RaggedRight\arraybackslash}p{3.5cm} >{\Centering\arraybackslash}p{1.1cm}}
\rowcolor{brand}
\thead{Fase} & \thead{Contenido} & \thead{Hito} & \thead{Sem.} \\
F0 — Inicio & \bodystrut Alcance, arquitectura, infraestructura base & H0: arquitectura definida & 1 \\
\rowcolor{rowtint}
F1 — Acceso & \bodystrut Login institucional, roles, sesiones, avisos & H1: acceso funcional & 2 \\
F2 — Búsqueda & \bodystrut Catálogo, búsqueda híbrida, filtros, paginación & H2: búsqueda de punta a punta & 4 \\
\rowcolor{rowtint}
F3 — Publicación & \bodystrut Borradores, carga, extracción/OCR, sugerencias IA, publicación & H3: publicación completa & 5 \\
F4 — Detalle & \bodystrut Página de lectura, descargas, relacionados & H4: detalle + relacionados & 6 \\
\rowcolor{rowtint}
F5 — Colaboración & \bodystrut Coautores, versiones de archivos & H5: colaboración + versiones & 7 \\
F6 — Ciclo de vida & \bodystrut Borrado/papelera/purga, reporte y moderación & H6: ciclo de vida completo & 8 \\
\rowcolor{rowtint}
F7 — Descubrimiento & \bodystrut Lecturas, “más leídos”, portada & H7: portada + descubrimiento & 8 \\
F8 — Despliegue & \bodystrut Nube, servicios de IA, salida a producción & H8: producción & 9 \\
\bottomrule[\heavyrulewidth]
\end{tabularx}

\end{document}
